The hypotheses of the next two corollaries on $b$ and on the maximizing sequence hold whenever $p\in\mathcal D_\omega^0$, since this gives $b>0$ and uniform integrability over the whole of $\mathscr G_\omega$ (\cref{thm:integrable-current-attainment}).

\Needspace{12\baselineskip}
\begin{corollary}\label{cor:semipositive-divisor-exclusion}
Assume $b>0$ and that $D_p$ has a smooth maximizing sequence $(T_j)$ with uniformly integrable densities $p k_{T_j}$. Let $E$ be an effective real divisor whose class has a smooth semipositive representative $\beta_E$. If
\[
 \int_{\{\beta_E=0\}}f\,\mathrm{d}V=0,
\]
then no threshold subsolution $u$ for $f$ satisfies $\omega+\ddc u\geq[E]$.
\end{corollary}
\begin{proof}
Suppose such a subsolution exists. By \cref{prop:threshold-divisorial-constraint}, there is a maximizing current $T\in\mathscr G_\omega$ with $\int_X\beta_E\wedge T=0$. We use the matrix-density notation of \cref{sec:weak-gauduchon-duality}: $\sigma_T=\omega\wedge T$ is the trace measure of $T$, $P_T$ is its matrix density and $B$ is the coefficient matrix of $\beta_E$. Since both matrices are positive semidefinite, the vanishing pairing gives $\operatorname{tr}(BP_T)=0$ at $\sigma_T$-almost every point. Wherever $B\ne0$, this forces $P_T$ to be singular: a positive definite $P_T$ would have strictly positive trace against the nonzero positive semidefinite matrix $B$. The determinant measure $\Gamma_0(T)$ vanishes where $P_T$ is singular, so it is concentrated on $\{\beta_E=0\}$. By \cref{thm:gaud-determinant-subcritical}, $\Gamma_0(T)=k_T\,\mathrm{d}V$, and the hypothesis on $f$ gives $p k_T=0$ almost everywhere. This contradicts $D_p(T)=\rho_\omega(p)>0$.
\end{proof}
